\begin{abstract}
	
	Quantum game theory extends classical strategic games by allowing players to apply unitary transformations to entangled quantum states. Within the Eisert–Wilkens–Lewenstein framework, equilibrium behavior is influenced by both entanglement and the admissible strategy space.
	
	This work interprets quantum equilibrium search as a coupled optimization problem over continuous parameterized unitary strategy spaces. From this perspective, strategic adaptation resembles reinforcement learning and continuous policy optimization over probabilistic payoff landscapes.
	
	We analyze the relationship between entanglement, equilibrium stability, admissible strategy classes, and optimization structure, while discussing limitations arising from nonconvexity, expanded strategy spaces, and physical noise.
	
\end{abstract}